\subsection{The codimension-three component for \texorpdfstring{$X_9$}{X9}}
\label{ssec:x9global}

We retain the toric fourfold $P$, the full coefficient space $V$ and the
exact four-node locus $S_4$ of Section~\ref{ssec:x9components}. Thus zero
coefficients are allowed, but additional singularities are excluded.
Number the lattice points as in \eqref{eq:x9coefficients}, and put
\begin{equation}\label{eq:x9determinantal}
 \mathcal D=\left\{c\in\C^{12}:\operatorname{rank}
 \begin{pmatrix}c_2&c_{10}&c_5&c_6\\c_7&c_8&c_4&c_3\end{pmatrix}
 \leq1\right\}.
\end{equation}
This irreducible determinantal variety has dimension nine. The special
form of $\Delta_9$ permits a classification of the part of $S_4$ carrying
a relation, as well as a projective small resolution on that entire part.

\begin{theorem}\label{thm:x9global}
The locus $S_4$ has exactly one connected component of codimension three,
namely $\mathcal D\cap S_4$. It meets the near-limit nodal construction of
Theorem~\ref{thm:main}. After marking and orienting the nodes by transport
from that construction, the integral relation lattice and the preimage of
ambient torsion are both $\Z(1,-1,-1,1)$.
Every other connected component of $S_4$ has codimension four and has
zero relation lattice and zero preimage of torsion. Every member of
$\mathcal D\cap S_4$ admits a projective small resolution.
\end{theorem}

The theorem does not count the codimension-four components;
Proposition~\ref{prop:x9different} shows that at least one exists. Its proof
uses the compact toric boundary, an evaluation-rank argument, and a
surface through the four nodes. The finite fan comparisons below use the
triangulation in Appendix~\ref{app:data-x9} and its accompanying data.\footnote{The
program \texttt{examples/verify\_x9\_global.py} reproduces the boundary
supports, singular ideals, surface fan and polytope comparisons used in
Sections~\ref{ssec:x9global}--\ref{ssec:x9reverse}.}

\subsubsection*{The compact boundary and a common affine chart}

Write $r_j$ for the primitive ray generators of the fan of $P$, in the
order of the accompanying triangulation. For a cone $\tau$, in addition
to $E_0(\tau)$ from Section~\ref{ssec:x9components}, define
\[
 E_{1,j}(\tau)=\{i:\langle n_i,r_j\rangle+1=1,
       \ \langle n_i,r_k\rangle+1=0\text{ for }k\in\tau,\ k\ne j\}.
\]
These are the supports of the first normal derivatives along the orbit
of $\tau$. Smoothness of the maximal cones makes this an ordinary
coordinate-derivative calculation.

\begin{lemma}\label{lem:x9boundaryglobal}
For every $c\in S_4$, the coefficients $c_0,c_1,c_3,c_4,c_5,c_6,c_7$
are nonzero. Every singular point lies in
\[
 U_{\langle r_8,r_{11}\rangle}
       =\C_x\times\C_y\times(\C^*)^2_{z,w},
 \qquad r_8=(1,-1,-1,-1),\quad r_{11}=(0,1,-1,-1),
\]
where the local equation is
\begin{equation}\label{eq:x9generalF}
\begin{split}
 F={}&c_1y^2+q(x)y+c_2x+c_3z+c_4w+c_5/z+c_6/w
                    +c_7xzw+c_8zw+c_{10},\\
 q(x)={}&c_0x^2+c_{11}x+c_9.
\end{split}
\end{equation}
\end{lemma}

\begin{proof}
The character change \eqref{eq:x9coordinates} is unimodular, and its
pairings with $r_8,r_{11}$ are the $x,y$ coordinate valuations. Multiplying
the Laurent equation by $xy$ gives \eqref{eq:x9generalF} on this chart.

The following cones have $E_0(\tau)=\{i\}$ and every $E_{1,j}(\tau)$
empty:
\[
\begin{array}{c|c|c}
i&\tau&\dim O(\tau)\\\hline
0&(19,120)&2\\
1&(136,141)&2\\
3&(112,134,155)&1\\
4&(121,144,149)&1\\
5&(127,149,150)&1\\
6&(8,132,157)&1
\end{array}
\]
If the indicated coefficient vanishes, the whole positive-dimensional
orbit is singular. This proves the first six nonvanishing assertions.

Suppose $c_7=0$. If $c_8=0$ also, the cone $(136,154,156)$ has
$E_0=\{7,8\}$ and all first-normal supports empty, again giving a
singular curve. If $c_8\ne0$, the four maximal cones
\[
 (9,80,114,136),\quad(10,85,117,136),\quad
 (114,131,136,154),\quad(117,131,136,156)
\]
give singular fixed points: each has $E_0=\{7\}$ and all
$E_{1,j}$ empty. There is a further singular point on each of the
following one-dimensional orbits:
\[
\begin{array}{c|c}
\tau&\text{only nonempty first-normal support}\\\hline
(9,114,154)&\{3,8\}\\
(10,117,156)&\{4,8\}\\
(9,80,136)&\{1,8\}\\
(10,85,136)&\{1,8\}
\end{array}
\]
Their orbit equations vanish identically, since $E_0=\{7\}$.
The remaining normal derivative is a binomial with nonzero coefficients
and distinct characters of $\C^*$, so has a zero on the orbit. At such a
zero every first derivative vanishes. These four orbits and the four
fixed points are disjoint, giving at least eight singular points.
Therefore $c_7\ne0$ on $S_4$.

Put $I=\{0,1,3,4,5,6,7\}$. For all $3,571$ nonzero cones outside
$\langle r_8,r_{11}\rangle$, the finite support calculation gives
\[
\begin{array}{l|r}
\text{an }i\in I\text{ outside the affine span of }
                  E_0(\tau)\setminus\{i\}&3,555\\
\text{otherwise, an }E_{1,j}(\tau)=\{i\},\quad i\in I&16.
\end{array}
\]
In the first row an affine-linear functional on the supported exponents
isolates the $i$th monomial. On the orbit it is a linear combination of
the equation and its logarithmic tangent derivatives, and cannot vanish
because $c_i\ne0$. In the second row a normal derivative is already a
nonzero monomial. Both arguments remain valid when any of the other
coefficients vanish. Together with the three nonzero cones $(8),(11)$
and $(8,11)$ in the stated chart, these exhaust the $3,574$ boundary
orbits. The dense torus is in the same chart.
\end{proof}

\subsubsection*{The coefficient-smoothing rank}

At a node in this chart the coefficient directions are evaluations of
\begin{equation}\label{eq:x9evalfunctions}
 x^2y,\ y^2,\ x,\ z,\ w,\ z^{-1},\ w^{-1},\ xzw,\ zw,\ y,\ 1,\ xy.
\end{equation}
These are directions in the full coefficient space even when a
coefficient of the defining equation is zero. Multiplying a local
equation by a unit only rescales the corresponding evaluation row.

\begin{lemma}\label{lem:x9rankglobal}
For $c\in S_4$, the coefficient-smoothing map has rank three if and
only if $c\in\mathcal D$. It never has rank at most two.
\end{lemma}

\begin{proof}
First, every three distinct singular points have independent evaluation
rows. If not, choose a minimally supported relation among at most three
rows. Two rows cannot be proportional, since $1,x,y,z,w$ distinguish
points, so all three weights are nonzero. The functions $1,y,y^2$ force
all three $y$-values to agree: three distinct values give a Vandermonde
matrix, and two values leave a singleton with zero weight. If their
common value is nonzero, $1,x,x^2y$ force common $x$ by the same argument.
If it is zero, $F_y=q(x)=0$ gives at most two $x$-values, and $1,x$ again
force common $x$. Now
\begin{equation}\label{eq:x9criticalpair}
 F_x=(2c_0x+c_{11})y+c_2+c_7zw=0,\qquad
 c_3z^2+(c_7x+c_8)(zw)z-c_5=0
\end{equation}
fix $p=zw$ and allow at most two $z$-values. Since $w=p/z$, three
distinct points are impossible. Lemma~\ref{lem:x9boundaryglobal}
supplies all nonzero coefficients used here.

Suppose now that the four rows have rank three, and let $\lambda$ be
their relation. It has full support by the preceding paragraph. Let
$W\subset\C^4$ be the space of evaluation vectors of the functions
in \eqref{eq:x9evalfunctions}. It has dimension three.

If the four $y$-values are distinct, $W=\langle1,y,y^2\rangle$.
The inclusions $x,xy\in W$, by interpolation at these four values,
give $x=\alpha+\beta y$. Then $x^2y\in W$ forces $\beta^2=0$.
Thus $x$ is constant, whereas $F_y=2c_1y+q(x)=0$ forces $y$ constant
as well, a contradiction. Three distinct $y$-values are also impossible:
the equations for $1,y,y^2$ make the sum of weights in each equal-value
group zero, and at least one group is a singleton.

If there are two $y$-values, the points form two pairs with opposite
weights within each pair. The equations for $x,xy$ force common $x$
within each pair. The two $x$-values differ, because $F_y=0$ determines
$y$ from $x$. Equation~\eqref{eq:x9criticalpair} fixes $p=zw$ within
each pair; its two $z$-values differ. Write $\Delta z_a,\Delta z_b$
for their differences, $S_a,S_b$ for their products, and
$\lambda_a,\lambda_b$ for the opposite weights. The functions $z,z^{-1}$
give
\[
 \lambda_a\Delta z_a+\lambda_b\Delta z_b=0,\qquad
 \frac{\lambda_a\Delta z_a}{S_a}
       +\frac{\lambda_b\Delta z_b}{S_b}=0.
\]
Both summands are nonzero, so $S_a=S_b$. Since
$\Delta w=-p\Delta z/S$, the equation for $w$ implies a common
product $p$ for both pairs. At each point, $zF_z-wF_w=0$ says
\begin{equation}\label{eq:x9cancel}
 (c_3+c_6/p)z^2=c_5+c_4p.
\end{equation}
If the left coefficient is nonzero, the two $z$-values in either pair
are opposite. The quadratic in \eqref{eq:x9criticalpair} then forces
$c_7x+c_8=0$ for both pairs, contradicting their distinct $x$-values.
Consequently both sides' coefficients in \eqref{eq:x9cancel} vanish.
On $zw=p$, the remaining equation is
\[
 h(x,y)=c_1y^2+q(x)y+Lx+M=0,\qquad
 L=c_2+c_7p,\quad M=c_{10}+c_8p.
\]
We have $h=h_x=h_y=0$ at two distinct $x$-values. Therefore
$q(x)^2-4c_1(Lx+M)$ has two distinct double roots and is the square of
a quadratic $s(x)$ with leading coefficient $c_0$. But
\[
 (q-s)(q+s)=4c_1(Lx+M)
\]
has right side of degree at most one and $\deg(q+s)=2$. Hence $q=s$
and $L=M=0$. Using both $h=0$ and $h_y=0$ now gives $y=0$, a
contradiction to two $y$-values.

There remains a common value $y_0$. If $y_0\ne0$, $F_y=0$ and the
quadratic in \eqref{eq:x9criticalpair} put the points in two pairs with
distinct $x$-values. The equations for $1,x$ give opposite weights in
each pair. The same calculation with $z,z^{-1},w$ gives common $p$.
Subtracting the two equations $F_x=0$ yields
$2c_0y_0(x_a-x_b)=0$, again impossible. Thus all four points have $y=0$.

Now $q(x)=0$ and $p=-c_2/c_7$ is fixed and nonzero. There must be two
distinct $x$-values, each with two distinct $z$-values. Equation
\eqref{eq:x9cancel} and the quadratic in \eqref{eq:x9criticalpair}
again force $c_3+c_6/p=c_5+c_4p=0$; the equation $F=0$ gives
$c_{10}+c_8p=0$. With $r=-1/p$ this is precisely
\[
 c_7=rc_2,\quad c_8=rc_{10},\quad c_4=rc_5,\quad c_3=rc_6,
\]
so $c\in\mathcal D$.

Conversely, at $c\in\mathcal D\cap S_4$, the nonvanishing of
$c_7,c_4,c_5$ forces $c_2\ne0$. Put
\begin{equation}\label{eq:x9factorparameters}
 A=c_2,\quad B=c_{10},\quad C=c_5,\quad D_0=c_6,\quad
 r=c_4/c_5.
\end{equation}
Then
\begin{equation}\label{eq:x9factorF}
 F=c_1y^2+q(x)y+G H,\qquad
 G=1+rzw,\quad H=Ax+B+C/z+D_0/w.
\end{equation}
The points $y=q=G=H=0$ are singular. Above each root of $q$, they are
given by
\begin{equation}\label{eq:x9fibrequadratic}
 -rD_0z^2+(Ax+B)z+C=0,\qquad w=-1/(rz).
\end{equation}
The leading and constant coefficients are nonzero. A repeated root of
$q$, or of \eqref{eq:x9fibrequadratic} above it, gives an actual singular
point with degenerate Hessian: at that point the differentials of
$y,q,G,H$ are dependent. Such a point is excluded by $S_4$. Thus there
are four distinct points, and
\begin{equation}\label{eq:x9nodejacobian}
 \det\frac{\partial(q,G,H)}{\partial(x,z,w)}
       =q'(x)r(C/z-D_0/w)\ne0.
\end{equation}
They are nodes by \eqref{eq:x9factorF}, and exhaust the singularities.
This includes a simple root $x=0$ when $c_9=0$.
At the nodes the evaluation functions span $\langle1,x,z,w\rangle$,
with relation $rD_0z+rCw=Ax+B$. The functions $1,x,z$ are independent
because $x$ takes two values and $z$ takes two values above each. Their
rank is exactly three.
\end{proof}

\subsubsection*{A compact surface and its small resolution}

\begin{lemma}\label{lem:x9surface}
For every $c\in\mathcal D\cap S_4$, the closure $S_r$ of
$\{y=0,\ 1+rzw=0\}$ is a smooth compact toric surface in $Y_f$.
Blowing up its ideal gives a projective small resolution of $Y_f$.
The surface has five toric boundary curves and $K_{S_r}^2=7$.
\end{lemma}

\begin{proof}
Set $R=r_{11}$ and $m=(4,2,1,1)$. The character $\chi^m=zw$ on the
toric divisor $E_R$ extends to a morphism $E_R\to\mathbb P^1$.
Indeed its pairings with the rays of each of the $18$ maximal cones of
the star of $R$ have one weak sign. The quotient lattice has coordinates
\[
 M/\Z R\longrightarrow\Z^3,\qquad
 v\longmapsto(v_0,v_1+v_2,v_1+v_3).
\]
The induced character is $(4,1,1)$, so its saturated kernel has
coordinates $(a,b)\mapsto(a,b,-4a-b)$. In this basis the fibre fan,
in positive cyclic order, is
\[
\begin{array}{c|rrrrr}
\text{ray index}&5&8&158&6&131\\\hline
\text{fibre ray}&(-1,1)&(1,-2)&(0,1)&(-1,3)&(-1,2).
\end{array}
\]
Its maximal cones are consecutive pairs, each of determinant one.
The self-intersections are $(0,1,-1,-1,-2)$, giving $K^2=12-5=7$.
The fibre over $-1/r\in\C^*$ is this smooth projective toric surface;
primitivity of the character makes its kernel torus connected. Fibres
over nonzero finite values are its torus translates.

The factorization \eqref{eq:x9factorF} shows that this fibre lies in
$Y_f$. It is $S_r$, and contains all four nodes. At a node, use
\eqref{eq:x9nodejacobian} and put $U=y$, $V=-(q+c_1y)$. The local
equation and centre become
\[
 UV=GH,\qquad S_r=\{U=G=0\}.
\]
The blowup is the incidence variety
\[
 U\beta=G\alpha,\qquad H\beta=V\alpha,
                  \qquad[\alpha:\beta]\in\mathbb P^1.
\]
On $\beta=1$ its free coordinates are $(G,V,\alpha)$, and on
$\alpha=1$ they are $(U,H,\beta)$. Thus it is smooth and replaces the
node by a $\mathbb P^1$ with normal bundle
$\mathcal O(-1)\oplus\mathcal O(-1)$. Away from the nodes, $S_r$ is a
Weil divisor in a smooth threefold and hence Cartier, so the blowup is
an isomorphism. The ideal blowup is projective over projective $Y_f$;
therefore the small resolution is projective.

Smallness and adjunction give trivial canonical bundle. The
anticanonical restriction sequence and toric vanishing give
$H^1(Y_f,\mathcal O)=H^2(Y_f,\mathcal O)=0$. The local small
resolution has vanishing higher direct images of $\mathcal O$:
the symmetric powers of the dual of its exceptional normal bundle are
sums of $\mathcal O(k)$, $k\ge0$, whose $H^1$ vanishes, and formal
functions applies. Leray gives the same cohomology vanishings on the
resolution. It is consequently a smooth projective Calabi--Yau threefold.
\end{proof}

\begin{proof}[Proof of Theorem~\ref{thm:x9global}]
We first exhibit a nonempty algebraic open subset of $\mathcal D$
contained in $S_4$. In \eqref{eq:x9factorparameters} take all boundary
coefficients equal to one and $c_{11}=0$. Its equation is
\[
 F=y^2+(x^2+1)y+(1+zw)(x+1+z^{-1}+w^{-1}).
\]
In $\Q[x,y,z,w,u]$, put $P_*=zwF$. The saturated singularity calculation is
\begin{equation}\label{eq:x9onenodeideal}
 (P_*,(P_*)_x,(P_*)_y,(P_*)_z,(P_*)_w,uzw-1)
   =(y,x^2+1,z+w-x-1,zw+1,uzw-1).
\end{equation}
The ideal is radical of degree four, and adjoining
$\det\operatorname{Hess}P_*$ gives the unit ideal. These facts also
follow by substituting $w=-1/z$, solving $x^2+1=0$ and
$z^2-(x+1)z-1=0$, and using \eqref{eq:x9nodejacobian}.
The boundary calculation in Lemma~\ref{lem:x9boundaryglobal} excludes
all other compact singularities for these nonzero essential
coefficients. Hence this member belongs to $S_4$.

On the parameter open where $r,A,C,D_0,c_0,c_1,c_9$ are nonzero,
$q$ has simple roots, and
\[
 \operatorname{Res}_x\bigl(q(x),(Ax+B)^2+4rCD_0\bigr)\ne0,
\]
the four expected nodes form a finite \'{e}tale degree-four scheme.
The nondegenerate Hessian makes it an open part of the relative singular
scheme; its finiteness makes it closed. The complementary singular
scheme is proper over the parameter space, as a closed subset of the
proper universal compact hypersurface. Removing its closed image gives
a Zariski-open subset $\mathcal D^{\rm good}\subset\mathcal D$ whose
members have precisely these four nodes. A finite singular locus also
implies reducedness. Equation~\eqref{eq:x9onenodeideal} proves that
this open subset is nonempty. It is connected and Euclidean dense in
the irreducible variety $\mathcal D$.

By Lemmas~\ref{lem:component-rank}, \ref{lem:component-smoothness} and
\ref{lem:x9rankglobal}, $\mathcal D^{\rm good}$ lies in a connected
component $C$ of $S_4$ of dimension nine, and is open in that component.
Every minor defining $\mathcal D$ vanishes on this open subset, hence
on $C$ by the identity theorem. Conversely, every point of
$\mathcal D\cap S_4$ is a limit of points of $\mathcal D^{\rm good}$
inside $S_4$; connected components are closed in $S_4$. Thus
$C=\mathcal D\cap S_4$. Lemma~\ref{lem:x9rankglobal} proves that all
other components have codimension four.

To meet the actual near-limit construction, choose $\mathscr C_0$
to contain the normalized coefficient vector $c^{(0)}$ with origin
coefficient $+1$ and all boundary coefficients $-1$. This vector lies
in $\mathcal Z$, since every parallelogram relation has two terms of
each sign. Use the height $H=k_0\varphi+h_C$ and coefficients
\[
 \gamma_i=-\tau^{H(n_i)}\quad(0\le i\le10),\qquad
 \gamma_{11}=1,\qquad0<\tau\ll1.
\]
These correspond to $c^{(0)}$. Linearity of $H$ on the four node
parallelograms gives
\[
 H(n_7)-H(n_2)=H(n_8)-H(n_{10})
             =H(n_4)-H(n_5)=H(n_3)-H(n_6).
\]
Thus $\gamma\in\mathcal D$, with $r>0$ and $A,B,C,D_0<0$.
The quadratic $q(x)=\gamma_0x^2+x+\gamma_9$ has two distinct positive
roots for small $\tau$, since all boundary heights are at least one.
Above each, the discriminant $(Ax+B)^2+4rCD_0$ is positive and
\eqref{eq:x9fibrequadratic} has two distinct nonzero roots.
These give four nodes. Lemma~\ref{lem:U}(1) puts $c^{(0)}$ in the
actual neighbourhood $U\subset U^+_\C$. Lemma~\ref{lem:W}(2) says
all compact singularities are among the four
unique critical points $q_P$. The four exhibited nodes therefore
exhaust them, so all critical values vanish and the member lies in the
actual $D_Q$ of Section~\ref{ssec:nodal}. This proves that $C$ is reached.

Lemma~\ref{lem:x9surface} gives the projective resolutions. It also
provides the integral relation directly. Remove small node balls from
the compact oriented surface $S_r$. Each boundary sphere is a primitive
generator of the third homology of the node link: it is the circle bundle
$\mathcal O(-1,-1)$ restricted to a ruling of the exceptional quadric.
The punctured surface and Lemma~\ref{lem:component-topology} therefore
give the relation $v=(1,1,1,1)$ after choosing these sphere orientations.
The rational kernel has rank one, so for the sphere map $d$,
\[
 \Z v\subseteq\ker d\subseteq d^{-1}(\Tors H_3)
       \subseteq\Q v\cap\Z^4=\Z v.
\]
All groups in this chain are equal. Proposition~\ref{prop:componentwise}
identifies the signs with $(1,-1,-1,1)$ using the marking transported
from the near-limit construction. On every other component the rational
sphere map has rank four, so both its integral kernel and its preimage
of torsion are zero.
\end{proof}
